\chapter{Abstract}
\begin{SingleSpace}

\textbf{Goal and object}. Ensuring security posture of critical IT infrastructures against insider attacks and threat actors by transitioning from reactive to preventive analysis of internal software components and introducing partial-automation to this secure lifecycle.\\\\
\textbf{Objectives and subjects}:
\begin{itemize}
    \item Constructed a conceptual framework for structured analysis of software in specialized isolated environments based on identified trust principles in the supply chain of IT products. It acts as a technical specification and standard foundation with a microkernel architecture, so such system would have reduced lines of code for its core logic, whereas third-party vendors should provide their modifications and implementations of components. Moreover, we demonstrated its potential integration into package distributions and software ecosystems as quality gates during the review stage, ensuring compliance with a predefined set of criteria.
    \item Applied a methodology of IaC (infrastructure-as-code) using definition files (combining shell scripts, scripting languages, and metadata) as both descriptive and execution artifacts by leveraging a graph-based execution engine and toolkits with such capabilities: portability, minimal external dependencies, and reliability.
    \item Conducted experiments utilizing an orchestrated ensemble of AI (artificial intelligence) agents based on self-hosted multimodal generative models to perform different tasks in provisioned isolated environments. Although such models are unreliable for regular tasks even with guardrail mechanisms, their fuzzy nature is applicable for automated testing and software pentests due to their ability to produce multiple solutions in complex cases and analyze large clusters of unstructured data. For example, a security analyst can spawn a large number of agents to rapidly iterate through various attack vectors and their combinations by generating small PoC (proof-of-concept) scripts/payloads, thereby operating more efficiently via adaptive automation. Moreover, it is possible to establish a full-cycle dynamic simulation that mimics the near-natural activity of autonomous users or developers based on predefined scenarios, behaviors, and goals. So, this method generates structured feedback and software evaluations, rather than simply stress-testing services deterministically like traditional fuzzers.
\end{itemize}
\textbf{Relevance and novelty}. Providing security measures and conducting attacks are both constant and coexisting processes that persist as long as there are numerous active actors with diverse intents. In these conditions, a preventive security approach can mitigate incidents at early stages and reduce potential risks by identifying vulnerabilities in software components. On the other hand, integration of AI agents and associated utilities remains relevant due to advancements of models in different domains. They serve as assistants for specialists and provide a natural language interface for machine interaction by automating multiple tasks such as autocompletion, boilerplate and script generation, translation, summarization, and automated code review. Also, the IT industry tries to shift toward optimized solutions and restricts the exposure of private data via small models on edge devices.\\\\
\textbf{Research methods}. Analyzing existing approaches based on industry-wide literature review in order to identify current gaps, potential solutions, and narrow the research scope. Furthermore, we applied a synthesis to develop new methodologies and derive experimental results.

\end{SingleSpace}
